\section{Technische und gestalterische Bewertung}
Die entwickelte WordClock kombiniert technische Funktionalität mit gestalterischen Anforderungen.  
Insbesondere die Verbindung aus sprachbasierter Zeitdarstellung, Zusatzinformationen und netzwerkbasierter Kommunikation stellte hohe Anforderungen an die Systemintegration.  

Die gewählte Hardwarearchitektur erwies sich insgesamt als zweckmäßig.  
Die eingesetzten Komponenten konnten zuverlässig miteinander kombiniert werden.  
Besonders die Aufteilung der Anzeige auf mehrere \gls{led}-Stränge erleichterte die Strukturierung der verschiedenen Funktionsbereiche.  

Auch die modulare Softwarearchitektur erwies sich als vorteilhaft.  
Die Trennung einzelner Funktionsmodule vereinfachte sowohl die Implementierung als auch die Fehlersuche während der Entwicklungsphase.  
Zusätzlich konnte die Verarbeitung externer Datenquellen zuverlässig in das Gesamtsystem integriert werden.  

Während der Integrationstests zeigte sich eine kurzzeitig hohe Leistungsaufnahme während des Bootvorgangs, die eine Anpassung der ursprünglichen Netzteildimensionierung erforderlich machte.  
Dadurch wurde deutlich, dass insbesondere bei großen \gls{led}-Systemen Einschaltvorgänge und Lastspitzen frühzeitig berücksichtigt werden müssen.  

Die Integration des Webservers und damit die Möglichkeit zur netzwerkbasierten Steuerung erweitert den Nutzen des Systems und erleichtert die Anpassung der Anzeige im Vergleich zu den in Kapitel~2 betrachteten und kommerziell vertriebenen WordClocks.  

Im gestalterischen Bereich konnte eine gleichmäßige und gut erkennbare Darstellung der Anzeigeelemente erreicht werden.  
Die entwickelte Konstruktion ermöglicht eine homogene Lichtverteilung sowie eine klare Trennung aktiver und inaktiver Anzeigeelemente.  
Gleichzeitig stellte die hohe Anzahl an Informationen erhöhte Anforderungen an die Anordnung und Lesbarkeit der Anzeige.
Das zeigt sich in der größe der WordClock. 
Im Vergleich zu den On-Market WordClocks ist jedoch eine präzisere und umfangreichere Anzeige möglich. 

Insgesamt zeigt die Umsetzung, dass die Kombination aus klassischer WordClock-Darstellung und erweiterten Informationsfunktionen technisch realisierbar ist.  
Gleichzeitig verdeutlicht das Projekt die Herausforderungen bei der Integration zahlreicher Funktionen innerhalb eines begrenzten Bauraums und einer begrenzten Anzeigefläche.